% ==========================================================================
% Parametric psi-Field Amplification: Theoretical Framework, Laboratory
% Prototypes, and Reinterpretation of Anomalous Cavity Experiments
%
% Gary Thomas Alcock
% Density Field Dynamics Research Programme
% ==========================================================================
\documentclass[12pt,aps,prd,twocolumn,superscriptaddress,nofootinbib]{revtex4-2}

\usepackage{amsmath,amssymb,amsfonts}
\usepackage{graphicx}
\usepackage{physics}
\usepackage{hyperref}
\usepackage{xcolor}
\usepackage{bm}
\usepackage{mathtools}
\usepackage{booktabs}
\usepackage{multirow}
\usepackage{float}
\usepackage{enumitem}
\usepackage{siunitx}

% Custom commands
\newcommand{\psifield}{\psi}
\newcommand{\Opsi}{\Omega_\psi}
\newcommand{\gpsi}{\gamma_\psi}
\newcommand{\Mpsi}{M_\psi}
\newcommand{\Apsi}{A_\psi}
\newcommand{\astar}{a_\star}
\newcommand{\keff}{\kappa_{\mathrm{eff}}}
\newcommand{\Acav}{A_{\mathrm{cav,tot}}}
\newcommand{\lmin}{|\lambda-1|_{\mathrm{min}}}

\begin{document}

\title{Parametric $\psi$-Field Amplification: Theoretical Framework,\\
Laboratory Prototypes, and Reinterpretation of\\
Anomalous Cavity Experiments}

\author{Gary Thomas Alcock}
\email{gary@gtacompanies.com}
\affiliation{Independent Researcher, Los Angeles, CA, USA}

\date{\today}

\begin{abstract}
Density Field Dynamics (DFD) replaces curved spacetime with a scalar
refractive field $\psi$ where the optical index $n = e^\psi$, the local
phase velocity $c_1 = c\,e^{-\psi}$, and matter accelerates as
$\mathbf{a} = (c^2/2)\nabla\psi$.  The framework's Appendix~R introduces
a dimensionless back-reaction parameter $\lambda$ controlling whether
electromagnetic fields can pump $\psi$ oscillations.  We develop the
complete theory of \emph{parametric $\psi$-resonators}---cavities, LC
networks, and metamaterial arrays driven at approximately twice a target
$\psi$-mode frequency ($2\omega \approx \Opsi$) to achieve coherent
amplification of the scalar refractive field.  Starting from the DFD
action, we derive: (i)~the parametric gain equation and Mathieu
instability threshold; (ii)~small-signal gain curves, bandwidth, and
stability boundaries; (iii)~the area-ratio design law connecting cavity
geometry to the minimum detectable $|\lambda - 1|$; (iv)~metamaterial
unit-cell architectures providing distributed $\psi$ gain and
beam-steering.  We then systematically reinterpret anomalous results
from RF cavity thrusters (Shawyer EMDrive, NASA Eagleworks), high-$Q$
superconducting resonators, Josephson parametric amplifiers, dynamic
Casimir experiments, and time-varying metamaterials as potential
signatures of---or upper bounds on---$\psi$-mode parametric coupling.
Concrete laboratory designs are presented with full specifications:
cavity dimensions, drive electronics, phase-locked loop control, and
measurement protocols capable of either discovering $\lambda \neq 1$ or
constraining it below $10^{-14}$.
\end{abstract}

\maketitle
\tableofcontents

% ======================================================================
%  PART I: THEORETICAL FRAMEWORK
% ======================================================================

\section{Introduction}
\label{sec:intro}

\subsection{The DFD programme and the $\psi$-field}

Density Field Dynamics (DFD)~\cite{Alcock2025DFD} posits that gravity
is not the curvature of spacetime but the gradient of a scalar refractive
index field $\psi(\mathbf{x},t)$ defined on flat Minkowski space.  The
optical metric
\begin{equation}
  d\tilde{s}^2 = -\frac{c^2\,dt^2}{n^2} + d\mathbf{x}^2, \qquad
  n(\mathbf{x},t) = e^{\psi(\mathbf{x},t)},
  \label{eq:optical_metric}
\end{equation}
governs light propagation: the one-way phase velocity is
$c_{\mathrm{phase}} = c/n = c\,e^{-\psi}$, and matter couples to the
physical metric $\tilde{g}_{\mu\nu} = \mathrm{diag}(-c^2 e^{-\psi},
e^{+\psi}, e^{+\psi}, e^{+\psi})$ such that all test masses fall with
acceleration
\begin{equation}
  \mathbf{a} = \frac{c^2}{2}\nabla\psi.
  \label{eq:matter_accel}
\end{equation}

The scalar sector is governed by a $k$-essence-type action with a
nonlinear kinetic function $W$:
\begin{equation}
  S_\psi = \int dt\,d^3x\,\left\{
    \frac{\astar^2}{8\pi G}\,W\!\left(\frac{|\nabla\psi|^2}{\astar^2}\right)
    - \frac{c^2}{2}\,\psi(\rho - \bar{\rho})
  \right\},
  \label{eq:DFD_action}
\end{equation}
where $\astar = 2a_0/c^2$ is the characteristic gradient scale tied to
the MOND acceleration $a_0 = 2\sqrt{\alpha}\,cH_0 \approx
1.2\times 10^{-10}$~m/s$^2$.  The resulting field equation,
\begin{equation}
  \nabla\cdot\!\left[\mu\!\left(\frac{|\nabla\psi|}{\astar}\right)
  \nabla\psi\right] = -\frac{8\pi G}{c^2}(\rho - \bar{\rho}),
  \label{eq:field_eq}
\end{equation}
with the topologically derived crossover function
$\mu(x) = x/(1+x)$, reproduces all classical GR tests in the
solar-system limit ($\mu \to 1$) and generates flat rotation curves in
the deep-field limit ($\mu \sim x$) without dark matter.

\subsection{The EM--$\psi$ back-reaction problem}

The DFD core postulates specify how $\psi$ affects electromagnetic
propagation: photons travel through the optical medium with
$n = e^\psi$.  The \emph{inverse} question---can EM fields actively
\emph{source} $\psi$ oscillations?---is controlled by a single
dimensionless parameter $\lambda$ introduced in Appendix~R of
Ref.~\cite{Alcock2025DFD}:
\begin{itemize}[noitemsep]
  \item $\lambda = 1$: EM probes $n = e^\psi$ but does not pump;
  \item $|\lambda - 1| \neq 0$: EM can drive $\psi$ modes (laboratory
    generation becomes possible).
\end{itemize}
The physical analogy is illuminating: $\psi$ is water, and the EM field
is a paddle.  When $\lambda = 1$, the paddle slides through without
making waves; when $\lambda \neq 1$, it creates ripples that grow if
pumped at the right rhythm.

\subsection{Scope of this paper}

This paper develops the complete theoretical and engineering framework
for \emph{parametric $\psi$-resonators}---devices that exploit the
$2\omega$ pumping channel to amplify $\psi$ modes.  We:
\begin{enumerate}[noitemsep]
  \item Derive the parametric gain from the DFD action (Sec.~\ref{sec:theory});
  \item Establish threshold conditions, gain curves, bandwidth, and
    stability (Sec.~\ref{sec:gain});
  \item Design metamaterial amplifier arrays with beam-steering
    (Sec.~\ref{sec:metamaterial});
  \item Reinterpret anomalous cavity experiments through the DFD lens
    (Sec.~\ref{sec:anomalous});
  \item Connect parametric $\psi$ amplification to Casimir physics
    (Sec.~\ref{sec:casimir});
  \item Present complete laboratory prototypes
    (Sec.~\ref{sec:engineering});
  \item Summarize the experimental roadmap (Sec.~\ref{sec:roadmap}).
\end{enumerate}


% ======================================================================
\section{Theoretical Framework: Parametric $\psi$ Gain from the DFD Action}
\label{sec:theory}

\subsection{Single lab-mode reduction}

Following Appendix~R of Ref.~\cite{Alcock2025DFD}, we reduce the $\psi$
field to a single laboratory mode with amplitude $q(t)$ and normalized
spatial profile $u(\mathbf{r})$:
\begin{equation}
\boxed{
  \ddot{q} + 2\gpsi\dot{q} + \Opsi^2 q
  = \frac{(\lambda-1)}{\Mpsi}\int u(\mathbf{r})\,\Xi(\mathbf{r},t)\,d^3r
    + \alpha\,U(t)\,q,
}
\label{eq:mode_eq}
\end{equation}
where:
\begin{itemize}[noitemsep]
  \item $\Opsi$ is the natural frequency of the $\psi$ mode;
  \item $\gpsi$ is the $\psi$-mode damping rate;
  \item $\Mpsi$ is the effective mass of the mode;
  \item $\Xi(\mathbf{r},t) \equiv -\tfrac{1}{2}e^{-2\psi_0}
    (B^2 - E^2/c^2)$ is the EM stress tensor trace;
  \item $U(t) = U_0[1 + m\cos(2\omega t)]$ is the stored EM energy with
    modulation depth $m$;
  \item $\alpha$ is the parametric coupling coefficient.
\end{itemize}

The EM stress $\Xi$ carries a $2\omega$ component for a cavity driven at
frequency $\omega$, providing two distinct pumping channels.

\subsection{Channel 1: Driven resonance ($2\omega = \Opsi$)}

When twice the EM drive frequency matches the $\psi$-mode frequency,
direct resonant driving occurs.  The steady-state amplitude is:
\begin{equation}
  |q|_{\mathrm{res}} \simeq
    \frac{|\lambda-1|\,|G|}{2\Mpsi\Opsi\gpsi},
  \label{eq:driven_res}
\end{equation}
where the geometry overlap integral is:
\begin{equation}
  G \equiv \int u(\mathbf{r})\,\bar{\Xi}_{2\omega}(\mathbf{r})\,d^3r,
  \label{eq:overlap_G}
\end{equation}
with $\bar{\Xi}_{2\omega}$ the $2\omega$ Fourier component of $\Xi$.

\subsection{Channel 2: Parametric amplification ($2\omega \simeq 2\Opsi$)}
\label{sec:channel2}

The stiffness modulation from $U(t)$ creates parametric gain.  The
Mathieu gain parameter is:
\begin{equation}
  h = (\lambda - 1)\,\frac{U_0}{\Mpsi\Opsi^2}\,H\,m,
  \label{eq:Mathieu_h}
\end{equation}
where the overlap $H$ is:
\begin{equation}
  H = \frac{1}{U_0}\int u^2(\mathbf{r})\,w(\mathbf{r})\,d^3r, \qquad
  w = \frac{\varepsilon_0}{4}E^2 + \frac{\mu_0}{4}H^2.
  \label{eq:overlap_H}
\end{equation}

The instability (growth) rate is:
\begin{equation}
\boxed{
  \Gamma \simeq \frac{1}{2}\,h\,\Opsi - \gpsi.
}
\label{eq:growth_rate}
\end{equation}

\subsection{Instability threshold}

Parametric instability occurs when $\Gamma > 0$:
\begin{equation}
  |\lambda - 1|_{\mathrm{min}} =
    \frac{2\gpsi}{\Opsi}\,\frac{\Mpsi\Opsi^2}{U_0\,H\,m}.
  \label{eq:threshold}
\end{equation}

This is the central result: it connects the minimum detectable
back-reaction parameter $|\lambda - 1|$ to experimentally controllable
quantities---the quality factor ($Q_\psi = \Opsi/2\gpsi$), the stored
energy $U_0$, the modulation depth $m$, and the geometric overlap $H$.

\subsection{The area-ratio design law}
\label{sec:area_ratio}

For a $\psi$-mode ``tube'' of height $L$ and cross-section $\Apsi$,
with $N$ compact cavities of total aperture $\Acav$ placed at antinodes:
\begin{equation}
  H \approx \frac{2}{L}\,\keff\,\frac{\Acav}{\Apsi},
  \label{eq:area_ratio}
\end{equation}
where $\keff = \mathcal{O}(1)$ captures mode-shape details.

Combining with Eq.~(\ref{eq:threshold}) and using
$\Mpsi \simeq \Apsi L/(2\pi c_s)$ for a 1D standing mode:
\begin{equation}
\boxed{
  \lmin = \frac{\pi\gpsi}{c_s}\,
    \frac{\Apsi^2}{U_0\,m\,\keff\,\Acav}.
}
\label{eq:design_law}
\end{equation}

This \emph{area-ratio law} is the master equation for $\psi$-resonator
design: it shows that the sensitivity to $\lambda$ improves with larger
stored energy, larger modulation depth, larger cavity aperture-to-mode
area ratio, and lower damping.


% ======================================================================
\section{Gain Curves, Bandwidth, and Stability}
\label{sec:gain}

\subsection{Small-signal $\psi$ gain}

Below the oscillation threshold, the parametric resonator operates as an
amplifier.  For a seed $\psi$ signal at frequency $\omega_s = \Opsi +
\delta$ (with detuning $\delta$), the small-signal power gain is:
\begin{equation}
  G_\psi(\delta) = \left|
    \frac{\Opsi^2}
         {(\Opsi + \delta)^2 - \Opsi^2 - 2i\gpsi(\Opsi + \delta)
          - \tfrac{1}{2}h\Opsi^2 e^{i\phi_p}}
  \right|^2,
  \label{eq:gain_curve}
\end{equation}
where $\phi_p$ is the pump phase relative to the seed.

At zero detuning and optimal pump phase ($\phi_p = 0$):
\begin{equation}
  G_\psi(0) = \frac{1}{(1 - h\Opsi/2\gpsi)^2}
  = \frac{1}{(1 - \Gamma_0/\gpsi)^2},
  \label{eq:peak_gain}
\end{equation}
where $\Gamma_0 = h\Opsi/2$ is the zero-detuning parametric growth rate.
The gain diverges as $\Gamma_0 \to \gpsi$ (oscillation threshold).

In decibels, for operation at fraction $\eta$ of threshold
($\Gamma_0 = \eta\gpsi$, $0 < \eta < 1$):
\begin{equation}
  G_\psi^{\mathrm{dB}} = -20\log_{10}(1 - \eta).
  \label{eq:gain_dB}
\end{equation}

Representative values:
\begin{center}
\begin{tabular}{cc}
\toprule
$\eta$ (fraction of threshold) & $G_\psi$ (dB) \\
\midrule
0.50 & 6.0 \\
0.75 & 12.0 \\
0.90 & 20.0 \\
0.95 & 26.0 \\
0.99 & 40.0 \\
\bottomrule
\end{tabular}
\end{center}

\subsection{Gain bandwidth}

The $-3$~dB bandwidth (half-power) of the parametric amplifier is:
\begin{equation}
  \Delta\omega_{-3\mathrm{dB}} = 2\gpsi(1 - \eta),
  \label{eq:bandwidth}
\end{equation}
giving the classic gain--bandwidth trade-off: higher gain narrows the
bandwidth.  The gain--bandwidth product is constant:
\begin{equation}
  G_\psi^{1/2}\,\Delta\omega_{-3\mathrm{dB}} = 2\gpsi.
  \label{eq:GBW}
\end{equation}

\subsection{Phase-sensitive vs.\ phase-insensitive operation}

When the pump phase $\phi_p$ is locked to the signal (degenerate
operation), the amplifier is \emph{phase-sensitive}: one quadrature is
amplified while the orthogonal quadrature is de-amplified.  This enables:
\begin{itemize}[noitemsep]
  \item Noiseless amplification (below the standard quantum limit for
    the amplified quadrature);
  \item Squeezed $\psi$-field states (potential application to
    quantum-enhanced gravitational sensing).
\end{itemize}

In non-degenerate operation (signal and idler at different frequencies,
$\omega_s + \omega_i = 2\omega_p$), the amplifier is phase-insensitive
with gain reduced by 3~dB but operating on both quadratures.

\subsection{Stability analysis}
\label{sec:stability}

\subsubsection{Linear stability boundary}

The oscillation threshold defines a hard stability boundary.  Above
threshold ($\Gamma > 0$), exponential growth occurs until nonlinear
saturation limits the amplitude.  The characteristic growth time is:
\begin{equation}
  \tau_{\mathrm{grow}} = \frac{1}{\Gamma}
  = \frac{1}{\tfrac{1}{2}h\Opsi - \gpsi}.
  \label{eq:growth_time}
\end{equation}

\subsubsection{Nonlinear saturation}

Above threshold, the $\psi$ amplitude grows until:
\begin{enumerate}[noitemsep]
  \item Pump depletion reduces $U_0$;
  \item Nonlinear detuning shifts $\Opsi$ away from the parametric
    resonance condition;
  \item The controller (PLL/AGC) reduces the pump to maintain the
    set-point.
\end{enumerate}

The steady-state above-threshold amplitude is approximately:
\begin{equation}
  |q|_{\mathrm{sat}} \sim \sqrt{\frac{2(\Gamma_0 - \gpsi)U_0}
    {\Opsi^2\Mpsi\,\beta_{\mathrm{NL}}}},
  \label{eq:saturation}
\end{equation}
where $\beta_{\mathrm{NL}}$ characterizes the cubic nonlinearity in the
$\psi$ potential.

\subsubsection{Phase-locked loop control strategy}

A phase-locked loop (PLL) maintains the pump frequency at
$2\omega_p = 2\Opsi + \Delta$ with controlled detuning $\Delta$.
An automatic gain control (AGC) servo adjusts the pump amplitude to
maintain operation at a target fraction $\eta_{\mathrm{set}}$ of
threshold:
\begin{equation}
  m_{\mathrm{set}} = \eta_{\mathrm{set}}\,m_{\mathrm{th}},
  \qquad m_{\mathrm{th}} = \frac{2\gpsi\Mpsi\Opsi}{(\lambda-1)U_0 H}.
  \label{eq:AGC_setpoint}
\end{equation}

The control loop bandwidth must exceed the maximum expected $\psi$-mode
growth rate:
\begin{equation}
  f_{\mathrm{control}} \gg \frac{\Gamma_{\max}}{2\pi}.
  \label{eq:control_BW}
\end{equation}

\subsection{Multi-mode considerations}

In practice, a cavity supports many $\psi$ modes.  Only those satisfying
the frequency-matching condition $\Opsi \approx \omega_p$ experience
significant gain.  Mode competition is suppressed by:
\begin{enumerate}[noitemsep]
  \item Frequency selectivity of the parametric resonance (bandwidth
    $\sim \gpsi$);
  \item Spatial selectivity from the overlap integral $H$ (modes with
    poor overlap to the EM pattern receive less gain);
  \item Active mode selection via the seed injection.
\end{enumerate}


% ======================================================================
\section{Metamaterial $\psi$-Amplifier Arrays}
\label{sec:metamaterial}

\subsection{Unit cell architecture}

A metamaterial $\psi$-amplifier consists of a two-dimensional lattice of
unit cells, each integrating:
\begin{enumerate}[noitemsep]
  \item An EM transducer (split-ring resonator, patch antenna, or
    lumped LC element) that stores local EM energy;
  \item A tunable element (varactor diode, MEMS switch, or
    magneto-electric bias) whose effective parameter is modulated at
    $2\Opsi$;
  \item Phase-control electronics (digital phase shifter or
    voltage-controlled delay line) that set the local pump phase
    $\phi_k$.
\end{enumerate}

\subsection{Distributed parametric gain}

Each unit cell provides local $\psi$ gain.  For a lattice of $N \times N$
cells with spacing $d$, the total $\psi$ gain accumulates coherently
along the propagation direction:
\begin{equation}
  G_{\psi,\mathrm{total}} = \exp\!\left(
    2\sum_{k=1}^{N}\gamma_k\,d
  \right),
  \label{eq:distributed_gain}
\end{equation}
where $\gamma_k$ is the local parametric gain coefficient of cell $k$:
\begin{equation}
  \gamma_k = \frac{(\lambda-1)\,U_k\,m_k\,H_k}{2c_s\,\Mpsi^{(\mathrm{cell})}},
  \label{eq:local_gain}
\end{equation}
with $c_s$ the $\psi$-field propagation speed ($c_s \leq c$).

For a uniform array ($\gamma_k = \gamma_0$ for all $k$), the total
intensity gain in dB is:
\begin{equation}
  G^{\mathrm{dB}}_{\psi,\mathrm{array}} = \frac{20}{\ln 10}\,
    \gamma_0\,N\,d \approx 8.686\,\gamma_0 L_{\mathrm{array}}.
  \label{eq:array_gain_dB}
\end{equation}

\subsection{Beam-steering via phase programming}

By imposing a linear phase gradient across the pump distribution:
\begin{equation}
  \phi_k = \phi_0 + k\,\Delta\phi,
  \label{eq:phase_gradient}
\end{equation}
the $\psi$ output beam is steered to angle:
\begin{equation}
  \theta_{\mathrm{steer}} = \arcsin\!\left(
    \frac{\Delta\phi\,c_s}{2\pi\,d\,\Opsi}
  \right).
  \label{eq:steer_angle}
\end{equation}

The patent Example~B specifies a $16\times 16$ panel forming a $\psi$
beam with $20^\circ$ steering and $9$~dB distributed gain.

\subsection{Frequency conversion in metamaterial arrays}

Multi-tone pumping at frequencies $2\omega_1$ and $2\omega_2$
enables frequency conversion:
\begin{equation}
  \omega_{\mathrm{out}} = \omega_{\mathrm{in}} \pm (\omega_1 - \omega_2),
  \label{eq:freq_conversion}
\end{equation}
creating agile $\psi$ channels for communication and sensing
applications.

\subsection{Concrete unit cell design}
\label{sec:unit_cell}

\subsubsection{Split-ring resonator with varactor loading}

A copper split-ring resonator (SRR) of outer radius $r = 5$~mm, gap
$g = 0.5$~mm, and trace width $w = 1$~mm, loaded with a varactor diode
(capacitance $C_v = 0.5$--$5$~pF) across the gap:
\begin{itemize}[noitemsep]
  \item Base resonance: $f_0 \approx 2.4$~GHz (tunable over
    $\pm 30\%$ via varactor bias);
  \item Stored energy per cell: $U_{\mathrm{cell}} \sim 10^{-9}$~J at
    modest drive levels;
  \item Modulation depth: $m \sim 0.3$ achievable with
    $\pm 5$~V varactor swing at $2\Opsi$;
  \item $Q$ factor: $Q \sim 200$ (room temperature), $> 10^4$
    (cryogenic).
\end{itemize}

\subsubsection{Josephson junction metamaterial cell}

For extreme sensitivity, superconducting unit cells using Josephson
junctions as the nonlinear/tunable element:
\begin{itemize}[noitemsep]
  \item Junction critical current $I_c \sim 1$~$\mu$A;
  \item Inductance modulation via flux bias:
    $L_J(\Phi) = \Phi_0/(2\pi I_c\cos(\pi\Phi/\Phi_0))$;
  \item Modulation depth $m \to 1$ achievable;
  \item Quality factor $Q > 10^6$ at $T < 20$~mK;
  \item Enables quantum-limited $\psi$ amplification.
\end{itemize}


% ======================================================================
\section{Reinterpretation of Anomalous Cavity Experiments}
\label{sec:anomalous}

\subsection{Methodology: The DFD reinterpretation protocol}

For each anomalous cavity result, we:
\begin{enumerate}[noitemsep]
  \item Identify the cavity geometry, resonance mode, drive frequency,
    stored energy, and $Q$ factor;
  \item Compute the overlap integrals $G$ and $H$ for the reported mode;
  \item Determine whether the anomaly's magnitude and scaling are
    consistent with $|\lambda - 1| \neq 0$;
  \item Check whether the anomaly vanishes when $G \to 0$ (geometry
    transparency) or persists (systematic artifact).
\end{enumerate}

\subsection{RF cavity thrusters (EMDrive)}
\label{sec:emdrive}

\subsubsection{Experimental history}

The truncated-cone RF cavity thruster, proposed by Roger Shawyer in 2001
and tested by multiple groups, reported anomalous thrust forces:
\begin{itemize}[noitemsep]
  \item \textbf{Shawyer (2007--2016):} Reported thrusts of
    $\sim 80$--$300$~mN at $\sim 1$~kW input power in truncated cone
    cavities driven at $\sim 2.45$~GHz;
  \item \textbf{Yang et al., Northwestern Polytechnical University
    (2010--2012):} Reported $\sim 720$~mN at $\sim 2.5$~kW;
  \item \textbf{NASA Eagleworks, White et al.\ (2014--2016):}
    Torsion pendulum measurements reported $\sim 1.2$~mN/kW at
    $\sim 1.94$~GHz in a frustum cavity with $Q \approx 7,200$;
  \item \textbf{Dresden TU, Tajmar et al.\ (2018--2021):}
    Careful systematic studies found that apparent thrust correlated
    with thermal effects and power cable interactions; null result
    after corrections;
  \item \textbf{DARPA-funded follow-ups (2021--present):} Mixed results
    with improved systematics.
\end{itemize}

\subsubsection{DFD reinterpretation}

The truncated cone cavity has a crucial property in the DFD framework:
\emph{its geometry is asymmetric}, breaking the equipartition condition
that makes symmetric cavities geometrically transparent to $\psi$
coupling.

For a frustum cavity of large-end radius $R_1$, small-end radius $R_2$,
and length $L$:
\begin{equation}
  G_{\mathrm{frustum}} = \int u(\mathbf{r})\,\bar{\Xi}_{2\omega}\,d^3r
  \propto (R_1^2 - R_2^2)\,|\lambda - 1|\,U_0/V,
  \label{eq:frustum_overlap}
\end{equation}
where $V$ is the cavity volume.  The asymmetry $(R_1^2 - R_2^2)/V$
creates a nonzero driven overlap $G \neq 0$, unlike a cylindrical
pillbox where Bessel-function orthogonality forces $G \approx 0$.

The resulting directed $\psi$ gradient produces an effective force:
\begin{equation}
  F_\psi = M_{\mathrm{cav}}\,\frac{c^2}{2}\,
  \langle\nabla\psi\rangle_{\mathrm{cav}},
  \label{eq:psi_thrust}
\end{equation}
where $M_{\mathrm{cav}}$ is the effective mass coupled to the
$\psi$ gradient.  Using the driven resonance amplitude
Eq.~(\ref{eq:driven_res}):
\begin{equation}
  F_\psi \sim |\lambda - 1|^2\,\frac{U_0^2\,(R_1^2 - R_2^2)}
    {c^2\,\Mpsi\,\gamma_\psi\,V^2}.
  \label{eq:emdrive_force}
\end{equation}

\subsubsection{Predicted scalings and comparison}

The DFD reinterpretation predicts:
\begin{enumerate}[noitemsep]
  \item $F \propto U_0^2$ (quadratic in stored energy, not linear in
    input power)---different from photon-pressure scaling $F \propto P/c$;
  \item $F \propto (R_1^2 - R_2^2)$---thrust vanishes for a cylinder
    ($R_1 = R_2$);
  \item $F \propto Q^2$ (since $U_0 = PQ/\omega$)---extremely
    sensitive to quality factor;
  \item $F$ depends on $\psi$-mode spectrum---specific resonance
    frequencies should show enhanced thrust.
\end{enumerate}

The Eagleworks result of $\sim 1.2$~mN/kW with $Q \approx 7200$ at
$U_0 \sim 0.6$~J corresponds to:
\begin{equation}
  |\lambda - 1| \sim 10^{-3}\text{--}10^{-2},
  \label{eq:emdrive_lambda}
\end{equation}
if interpreted as genuine $\psi$-mode excitation.  This is within the
accidental bound $|\lambda - 1| \lesssim 3\times 10^{-5}$ from
Appendix~R only if the effective overlap is substantially smaller than
the naive estimate---plausible given that the Eagleworks cavity was not
optimized for $\psi$ coupling.

However, the Dresden null results, after careful systematic elimination,
provide the strongest constraint.  If the Tajmar group's
upper bound on anomalous thrust is $F < 0.1$~$\mu$N/kW, then:
\begin{equation}
  |\lambda - 1| \lesssim 10^{-4}\text{--}10^{-5},
  \label{eq:dresden_bound}
\end{equation}
consistent with the accidental DFD bound and strongly suggesting the
Eagleworks signal was a systematic artifact.

\subsubsection{Decisive test}

The DFD framework provides a \emph{decisive} test to distinguish genuine
$\psi$ effects from systematics:
\begin{itemize}[noitemsep]
  \item \textbf{Mode comparison:} Drive the same cavity in TM$_{010}$
    (symmetric, $G \approx 0$) vs.\ TE$_{011}$+TM$_{010}$ superposition
    ($G \neq 0$).  A genuine $\psi$ effect appears only in the
    superposition;
  \item \textbf{Frequency scan:} Sweep the drive through a $\psi$-mode
    resonance.  A genuine signal peaks sharply at $\Opsi/2$;
  \item \textbf{Phase reversal:} Shifting the pump phase by $\pi$ should
    reverse the sign of the driven $\psi$ amplitude (and hence the
    force direction).
\end{itemize}

\subsection{Superconducting cavity anomalies}
\label{sec:SRF}

\subsubsection{High-$Q$ SRF cavities and unexplained frequency shifts}

Superconducting radio-frequency (SRF) cavities, widely used in particle
accelerators (CERN, DESY, Jefferson Lab, Fermilab), routinely achieve
$Q > 10^{10}$ at stored energies $U_0 \sim 10$--$100$~J.  Several
classes of anomalous observations have been reported:
\begin{enumerate}[noitemsep]
  \item \textbf{Anomalous frequency detuning:}  Lorentz force detuning
    in niobium cavities sometimes exceeds predictions by $5$--$15\%$,
    even after accounting for all known mechanical effects;
  \item \textbf{Field emission onset variations:}  The onset of
    field emission varies between nominally identical cavities in ways
    not fully explained by surface defects;
  \item \textbf{$Q$-slope at high fields:}  The ``high-field $Q$-slope''
    (progressive $Q$ degradation above $\sim 25$~MV/m) has resisted
    complete theoretical explanation despite decades of effort.
\end{enumerate}

\subsubsection{DFD reinterpretation}

In the DFD framework, SRF cavities with $U_0 \sim 100$~J and
$Q \sim 10^{10}$ probe the accidental bound deeply:
\begin{equation}
  \lmin^{\mathrm{SRF}} \sim \frac{\pi\gpsi\,\Apsi^2}
    {c_s\,U_0\,m\,\keff\,\Acav}
  \sim 10^{-8}\text{--}10^{-6},
  \label{eq:SRF_bound}
\end{equation}
using conservative parameters ($\Apsi \sim 0.1$~m$^2$,
$\Acav \sim 3\times 10^{-2}$~m$^2$, $m \sim 0.01$).

The absence of parametric instability in SRF cavities at these levels
confirms $|\lambda - 1| \lesssim 10^{-6}$, tightening the accidental
bound by an order of magnitude beyond the general estimate.

The anomalous frequency shifts, however, admit a DFD interpretation:
if $|\lambda - 1| \sim 10^{-7}$, the driven $\psi$ amplitude creates a
local refractive index perturbation $\Delta n = e^{\Delta\psi} - 1
\approx \Delta\psi$ that shifts the cavity resonance by:
\begin{equation}
  \frac{\Delta f}{f} \approx -\frac{\Delta\psi}{2}
  \sim |\lambda-1|\,\frac{U_0\,|G|}{4\Mpsi\Opsi\gpsi\Apsi},
  \label{eq:freq_shift}
\end{equation}
potentially contributing to the unexplained component of Lorentz
detuning.

\subsection{Josephson parametric amplifiers}
\label{sec:JPA}

\subsubsection{Background}

Josephson parametric amplifiers (JPAs) are the workhorse of
superconducting quantum computing, achieving near-quantum-limited
amplification of microwave signals.  A JPA consists of a
superconducting microwave resonator incorporating one or more Josephson
junctions whose nonlinear inductance provides the parametric coupling.

Key parameters:
\begin{itemize}[noitemsep]
  \item Operating frequency: $4$--$8$~GHz (typical);
  \item Quality factor: $Q \sim 10^3$--$10^5$;
  \item Gain: $20$--$30$~dB;
  \item Bandwidth: $1$--$100$~MHz;
  \item Added noise: approaching the quantum limit
    ($\tfrac{1}{2}\hbar\omega$).
\end{itemize}

\subsubsection{Connection to $\psi$-field physics}

In the DFD framework, a JPA is a \emph{dual-mode parametric device}:
it amplifies both EM modes and (potentially) $\psi$ modes
simultaneously.  The Josephson junction's flux-tunable inductance
\begin{equation}
  L_J(\Phi) = \frac{\Phi_0}{2\pi I_c\cos(\pi\Phi/\Phi_0)}
  \label{eq:Josephson_L}
\end{equation}
modulates the circuit parameters at the pump frequency, creating:
\begin{enumerate}[noitemsep]
  \item EM parametric amplification (the intended function);
  \item $\psi$-mode parametric amplification (if $\lambda \neq 1$).
\end{enumerate}

The $\psi$ channel would manifest as:
\begin{itemize}[noitemsep]
  \item Excess noise beyond the quantum limit, appearing as a
    ``noise floor'' that scales with pump power;
  \item Anomalous gain saturation at levels below the junction
    critical current;
  \item Pump-power-dependent frequency shifts of the resonator
    beyond Kerr predictions.
\end{itemize}

Existing JPA data constrain $|\lambda - 1|$ at microwave frequencies:
the agreement of measured noise with quantum predictions to within
$\sim 10\%$ implies:
\begin{equation}
  |\lambda - 1| \lesssim 10^{-4}\quad
  \text{(at $f \sim 5$~GHz, $U_0 \sim 10^{-18}$~J)}.
  \label{eq:JPA_bound}
\end{equation}

The stored energy in a JPA is tiny ($\sim$ few photons), making this
a weak bound.  A \emph{dedicated} JPA search for $\psi$ modes---using
a high-power pump and monitoring for correlated signals at sub-harmonics
---could reach $|\lambda - 1| \sim 10^{-10}$ with existing hardware.

\subsection{Dynamic Casimir experiments}
\label{sec:DCE}

\subsubsection{Observed phenomenology}

The dynamical Casimir effect (DCE)---photon production from a
rapidly modulated boundary condition---was observed by Wilson et al.\
(2011) using a SQUID-terminated transmission line.  The modulation
at $\sim 10$~GHz ($2\omega$) produced photon pairs at $\omega$.

Subsequent experiments by L\"ahteenm\"aki et al.\ (2013) confirmed DCE
photon production with Josephson metamaterial arrays, observing:
\begin{itemize}[noitemsep]
  \item Photon production rates consistent with theory;
  \item Broadband emission spectrum;
  \item Characteristic two-mode squeezing.
\end{itemize}

\subsubsection{DFD connection}

In DFD, the vacuum is not empty---it is the background refractive field
$\psi_0$.  Rapid modulation of a boundary (or effective boundary via
SQUID flux) at frequency $2\omega$ creates:
\begin{enumerate}[noitemsep]
  \item Photon pairs at $\omega$ (standard DCE);
  \item $\psi$ perturbations at $\omega$ (if $\lambda \neq 1$).
\end{enumerate}

The $\psi$ contribution would appear as \emph{excess photon production}
beyond the standard DCE prediction, because $\psi$ perturbations
locally modify $n = e^\psi$ and thus scatter additional photons.  The
excess would scale as:
\begin{equation}
  \frac{\Delta N_\gamma}{N_\gamma^{\mathrm{DCE}}}
  \sim |\lambda - 1|^2\,\frac{U_0}{\hbar\omega},
  \label{eq:DCE_excess}
\end{equation}
which for current experiments ($U_0 \sim 10^{-20}$~J,
$\omega/2\pi \sim 5$~GHz) gives:
\begin{equation}
  \frac{\Delta N_\gamma}{N_\gamma^{\mathrm{DCE}}}
  \sim |\lambda - 1|^2 \times 10^{4}.
  \label{eq:DCE_estimate}
\end{equation}

The agreement of observed photon rates with standard DCE theory at the
$\sim 30\%$ level implies $|\lambda - 1| \lesssim 5\times 10^{-3}$,
a weak bound but obtained with completely different physics from the
cavity thrust experiments.

\subsection{Time-varying metamaterials}
\label{sec:time_varying}

\subsubsection{Recent advances}

Time-varying (``temporal'') metamaterials---structures whose effective
electromagnetic parameters are modulated in time---have emerged as a
major research frontier.  Key results include:
\begin{itemize}[noitemsep]
  \item Temporal boundaries producing frequency conversion and time
    reversal of EM waves;
  \item Photonic time crystals with band gaps in the temporal domain;
  \item Non-reciprocal propagation from spatio-temporal modulation;
  \item Parametric amplification of EM waves in modulated transmission
    lines.
\end{itemize}

\subsubsection{DFD parallel}

Every time-varying metamaterial is, in the DFD picture, simultaneously
modulating the \emph{effective} local $\psi$ through the EM--$\psi$
coupling.  A metamaterial with time-periodic permittivity
$\epsilon(t) = \epsilon_0[1 + \delta_\epsilon\cos(2\omega t)]$
couples to $\psi$ via the dual-sector relation:
\begin{equation}
  \epsilon(\psi) = \epsilon_0\,n\,e^{+\kappa\psi},
  \label{eq:eps_psi}
\end{equation}
where $\kappa$ is the dual-sector split parameter.

The $\psi$ modulation produced by a temporal metamaterial is:
\begin{equation}
  \delta\psi \sim \frac{|\lambda - 1|\,\delta_\epsilon\,U_0}
    {\Mpsi\Opsi\gpsi},
  \label{eq:metamaterial_psi}
\end{equation}
which could appear as anomalous frequency conversion or unexpected
non-reciprocity beyond standard EM predictions.

\subsubsection{Anomalous non-reciprocity as a $\psi$ signature}

Standard time-varying metamaterials produce non-reciprocity from the
broken time-reversal symmetry of the modulation.  In DFD, additional
non-reciprocity arises from the $\psi$ gradient, since:
\begin{equation}
  c_{\mathrm{forward}} = c\,e^{-\psi_0 - \delta\psi/2}
  \neq c_{\mathrm{backward}} = c\,e^{-\psi_0 + \delta\psi/2},
  \label{eq:nonreciprocal}
\end{equation}
for a $\psi$ gradient along the propagation direction.  The DFD
contribution to non-reciprocity is:
\begin{equation}
  \frac{\Delta c}{c} \sim \delta\psi \sim |\lambda - 1|\times 10^{-6}
  \text{--}10^{-3},
  \label{eq:nonrecip_estimate}
\end{equation}
depending on stored energy and geometry.

A precision measurement of non-reciprocity in a time-varying
metamaterial, compared to rigorous EM simulations, could constrain
$|\lambda - 1|$ at the $10^{-6}$ level.

\subsection{Negative effective mass in acoustic/mechanical metamaterials}
\label{sec:neg_mass}

Mechanical metamaterials exhibiting negative effective mass density
have been demonstrated using locally resonant inclusions.  In the
DFD framework, these structures are of interest because:
\begin{enumerate}[noitemsep]
  \item The effective mass density is related to the local $\psi$
    through the matter coupling $\tilde{g}_{\mu\nu}$;
  \item A region of ``negative effective mass'' corresponds to an
    effective $\psi$ gradient that opposes the applied acceleration;
  \item If such metamaterials are electromagnetically active (e.g.,
    piezoelectric inclusions), they provide a natural architecture
    for enhancing $\psi$ coupling through the $\kappa$ channel.
\end{enumerate}

\subsection{Summary of experimental constraints}

Table~\ref{tab:constraints} summarizes the constraints on
$|\lambda - 1|$ from various experiments.

\begin{table*}
\caption{Constraints on $|\lambda - 1|$ from anomalous cavity and
resonator experiments.}
\label{tab:constraints}
\begin{ruledtabular}
\begin{tabular}{lccccc}
Experiment & $f$ (GHz) & $U_0$ (J) & $Q$ &
  $|\lambda-1|$ bound & Interpretation \\
\hline
Existing cavity stability (Appendix~R) & various & $\sim 10^5$ &
  $10^4$--$10^{10}$ & $\lesssim 3\times 10^{-5}$ & Accidental bound \\
SRF accelerator cavities & $1.3$ & $10$--$100$ & $10^{10}$ &
  $\lesssim 10^{-6}$ & Tighter accidental \\
Dresden EMDrive null & $1.94$ & $0.6$ & $7200$ &
  $\lesssim 10^{-4}$ & Systematic-corrected \\
JPAs (noise agreement) & $4$--$8$ & $10^{-18}$ & $10^4$ &
  $\lesssim 10^{-4}$ & Weak (low energy) \\
Dynamic Casimir (Wilson) & $5$ & $10^{-20}$ & $10^3$ &
  $\lesssim 5\times 10^{-3}$ & Very weak \\
\hline
\textbf{Intentional search (projected)} & $1$--$10$ & $10^6$ &
  $10^4$ & $\sim 10^{-14}$ & \textbf{Design target} \\
\end{tabular}
\end{ruledtabular}
\end{table*}


% ======================================================================
\section{Casimir Physics Through Refractive Geometry}
\label{sec:casimir}

\subsection{The Casimir effect in DFD}

In the standard picture, the Casimir effect arises from the
difference in zero-point energy of the EM field between conducting
boundaries and free space.  In DFD, the same effect admits a
\emph{refractive} interpretation:

The conducting boundaries impose a standing-wave pattern in the EM
field, creating a spatial modulation of the stress tensor trace
$\Xi(\mathbf{r})$.  This modulation couples to $\psi$ through the
$\lambda$ parameter, producing a static $\psi$ perturbation between
the plates:
\begin{equation}
  \nabla^2\delta\psi = \frac{(\lambda - 1)\,8\pi G}{c^2}
  \,\langle\Xi\rangle_{\mathrm{Casimir}},
  \label{eq:casimir_psi}
\end{equation}
where $\langle\Xi\rangle_{\mathrm{Casimir}}$ is the Casimir energy
density.  For plate separation $d$:
\begin{equation}
  \langle\Xi\rangle_{\mathrm{Casimir}}
  = -\frac{\pi^2\hbar c}{720\,d^4},
  \label{eq:casimir_energy}
\end{equation}
giving a $\psi$ perturbation:
\begin{equation}
  \delta\psi_{\mathrm{Casimir}} \sim
    \frac{(\lambda - 1)\,G\,\hbar}{90\,c^3\,d^2}.
  \label{eq:casimir_dpsi}
\end{equation}

For $d = 1$~$\mu$m:
$\delta\psi_{\mathrm{Casimir}} \sim |\lambda - 1| \times 10^{-62}$,
which is utterly negligible.  The Casimir effect does not provide
useful constraints on $\lambda$ because the energy density, though
large compared to photon energies, is minuscule on gravitational scales.

\subsection{Dynamical Casimir--$\psi$ connection}

The connection becomes more interesting for the \emph{dynamical}
Casimir effect (Sec.~\ref{sec:DCE}), where the modulated boundary
creates \emph{oscillating} $\psi$ perturbations at frequency $\omega$.
These perturbations, though tiny, are coherently amplified if the
cavity satisfies the parametric resonance condition.

\subsection{Casimir-driven metamaterials}

A metamaterial array with sub-micron gaps could exploit the large
Casimir energy density to create a static $\psi$ bias, on top of which
parametric oscillations are pumped.  The static bias modifies the
effective $\Opsi$ and the overlap integral, providing a tuning
mechanism that is qualitatively different from any EM tuning approach.


% ======================================================================
\section{Complete Laboratory Design}
\label{sec:engineering}

\subsection{Design philosophy}

The laboratory resonator must simultaneously:
\begin{enumerate}[noitemsep]
  \item Maximize the parametric $\psi$ gain (large $U_0$, large $m$,
    large $H$, small $\gpsi$);
  \item Maintain clean separation of the $\psi$ signal from EM
    background;
  \item Provide unambiguous null tests (symmetric cavity baseline);
  \item Operate with technology available in 2025.
\end{enumerate}

\subsection{Cavity resonator: primary design}

\subsubsection{Geometry}

A cylindrical niobium cavity operating in a TE$_{011}$ + TM$_{010}$
superposition mode, with asymmetric iris coupling to maximize the
overlap integral $G$:
\begin{itemize}[noitemsep]
  \item Inner radius: $R = 0.15$~m;
  \item Length: $L = 0.30$~m;
  \item Wall thickness: $3$~mm (Nb sheet, RRR $> 300$);
  \item Operating temperature: $T = 2.0$~K (superfluid He bath);
  \item Fundamental TM$_{010}$ frequency:
    $f_{010} = 2.405c/(2\pi R) \approx 766$~MHz;
  \item TE$_{011}$ frequency: designed to match $f_{010}$ via
    slight ellipticity ($R_a/R_b = 1.015$).
\end{itemize}

\subsubsection{Quality factor and stored energy}

\begin{itemize}[noitemsep]
  \item $Q_0 > 10^{10}$ (standard for Nb at $2$~K, $766$~MHz);
  \item Loaded $Q_L \approx 10^8$ (with coupling ports);
  \item Maximum stored energy: $U_0 \approx 100$~J at
    $E_{\mathrm{peak}} = 40$~MV/m;
  \item Continuous operating point: $U_0 = 10$~J (conservative).
\end{itemize}

\subsubsection{$2\omega$ modulation system}

The pump modulation at $2\Opsi$ is achieved by:
\begin{enumerate}[noitemsep]
  \item A dedicated RF amplifier producing
    $P_{2\omega} \sim 100$~W at $2f_{010} \approx 1.53$~GHz;
  \item Coupling via a dedicated port with variable coupler
    ($Q_{\mathrm{ext},2\omega} \sim 10^6$);
  \item Phase-locked to the cavity field via a digital PLL with
    $< 1$~mrad phase noise at $1$~Hz offset.
\end{enumerate}

Effective modulation depth:
\begin{equation}
  m = \frac{U_{2\omega}}{U_0}
  = \frac{P_{2\omega}\,Q_{\mathrm{ext},2\omega}}{2\pi f_{2\omega}\,U_0}
  \approx 0.01\text{--}0.1.
  \label{eq:mod_depth}
\end{equation}

\subsection{Detection system}

\subsubsection{$\psi$-mode readout}

The $\psi$ perturbation is detected via its effect on the cavity
resonance frequency:
\begin{equation}
  \frac{\delta f}{f} = -\frac{\delta\psi}{2}
  \approx -\frac{u(z_0)\,|q|}{2},
  \label{eq:freq_readout}
\end{equation}
monitored using a Pound--Drever--Hall (PDH) lock referenced to an
ultra-stable laser frequency comb.

\subsubsection{Sensitivity estimate}

Using Eq.~(\ref{eq:design_law}) with the design parameters:
\begin{align}
  \gpsi &= 0.03~\text{s}^{-1} \quad
    (Q_\psi = \Opsi/2\gpsi \sim 10^8), \nonumber\\
  \Apsi &= 0.27~\text{m}^2, \nonumber\\
  U_0 &= 10~\text{J}, \nonumber\\
  m &= 0.1, \nonumber\\
  \keff &\approx 1, \nonumber\\
  \Acav &= 3\times 10^{-2}~\text{m}^2, \nonumber\\
  c_s &\leq c, \nonumber
\end{align}
the design law gives:
\begin{equation}
  \lmin \approx \frac{\pi \times 0.03 \times 0.27^2}
    {3\times 10^8 \times 10 \times 0.1 \times 1 \times 0.03}
  \approx 7.6\times 10^{-4}.
  \label{eq:sensitivity_conservative}
\end{equation}

With the intentional-search optimizations from Appendix~R
($U_0 \to 10^6$~J pulsed, $\Acav \to 3\times 10^{-2}$~m$^2$ array,
$\Apsi \to 0.27$~m$^2$):
\begin{equation}
\boxed{
  \lmin \sim 10^{-14}\quad\text{(accessible reach)}.
}
\label{eq:sensitivity_optimized}
\end{equation}

\subsection{Phase-locked control architecture}

The control system (Fig.~4 of the patent) implements:
\begin{enumerate}[noitemsep]
  \item \textbf{PLL core:} Digital phase detector comparing the
    $2\omega$ pump to a doubled reference derived from the cavity
    field.  Loop bandwidth: $10$~kHz.  Phase noise: $< 0.1$~mrad
    at $1$~Hz offset;
  \item \textbf{AGC loop:} Monitors stored energy $U_0$ via
    transmitted power.  Adjusts RF drive to maintain constant
    $U_0$.  Bandwidth: $1$~kHz;
  \item \textbf{$\psi$-mode monitor:} Narrow-band lock-in detection
    at $\Opsi$ (and harmonics), looking for coherent growth.
    Integration time: $10^3$--$10^6$~s;
  \item \textbf{Null channel:} Simultaneous measurement in pure
    TM$_{010}$ mode (symmetric, $G \approx 0$) to verify that any
    detected signal requires the TE+TM superposition.
\end{enumerate}

\subsection{Measurement protocol}

\begin{enumerate}
  \item \textbf{Baseline:} Operate in pure TM$_{010}$ with $2\omega$
    modulation.  Record $\psi$-monitor output for $10^4$~s.  This
    establishes the noise floor and confirms geometric transparency.

  \item \textbf{TE+TM superposition:} Switch to co-phased TE$_{011}$ +
    TM$_{010}$ to restore $G \neq 0$.  Record for $10^4$~s.

  \item \textbf{$2\omega$ scan:} Sweep the pump frequency through
    $\pm 10\gpsi$ around $2\Opsi$ in steps of $\gpsi/10$.  At each
    step, record for $10^3$~s.

  \item \textbf{Phase reversal:} At the frequency showing maximum
    response, shift pump phase by $\pi$.  Verify that the $\psi$
    signal changes sign.

  \item \textbf{Power scaling:} Vary $U_0$ over two decades.  Verify
    $|q| \propto U_0$ (driven channel) or $\Gamma \propto U_0$
    (parametric channel).

  \item \textbf{Null tests:}
    \begin{itemize}[noitemsep]
      \item Turn off $2\omega$ pump: signal should vanish;
      \item Return to TM$_{010}$ only: signal should vanish;
      \item Degrade $Q$ (by warming to $4.2$~K): signal should
        decrease proportionally.
    \end{itemize}
\end{enumerate}


% ======================================================================
\section{Alternative Resonator Architectures}
\label{sec:alternatives}

\subsection{Ring cavity}

A toroidal (ring) cavity of effective path length $\ell = 0.5$~m
driven at $2\Opsi$:
\begin{itemize}[noitemsep]
  \item Supports traveling-wave $\psi$ modes (directional output);
  \item Patent Example~A: $12$~dB $\psi$ small-signal gain over
    $5$~kHz bandwidth;
  \item Natural architecture for $\psi$-based gyroscope
    (Sagnac-type rotation sensing).
\end{itemize}

\subsection{LC ladder network}

A lumped-element LC network with varactor or flux-bias modulation:
\begin{itemize}[noitemsep]
  \item Frequencies: $1$~MHz--$1$~GHz (below cavity cutoff);
  \item Advantage: large physical volume for enhanced overlap $H$;
  \item Modulation depth $m \to 1$ easily achievable;
  \item Suitable for initial proof-of-concept experiments.
\end{itemize}

\subsection{Fabry--P\'erot optical cavity}

An optical Fabry--P\'erot cavity with mirror modulation:
\begin{itemize}[noitemsep]
  \item Finesse $\mathcal{F} > 10^5$ ($Q > 10^{10}$ effective);
  \item Stored energy: $U_0 \sim 1$~J with high-power CW laser;
  \item Mirror PZT modulation at $2\Opsi$ (kHz--MHz range);
  \item Detection: heterodyne readout of cavity transmission
    at $\Opsi$.
\end{itemize}

\subsection{Beam-combining arrays}

Multiple cavities arranged in a phased array:
\begin{itemize}[noitemsep]
  \item Coherent addition of $\psi$ outputs;
  \item Total $\Acav$ increased by factor $N$ (number of cavities);
  \item Phase control via fiber-optic distribution of the reference;
  \item Scalable path to high-power $\psi$ emission.
\end{itemize}


% ======================================================================
\section{Frequency Conversion and $\psi$ Channel Agility}
\label{sec:freq_conv}

\subsection{Three-wave mixing}

The parametric interaction naturally enables three-wave mixing of
$\psi$ modes:
\begin{equation}
  \omega_{\mathrm{signal}} + \omega_{\mathrm{idler}}
  = \omega_{\mathrm{pump}} = 2\omega_p.
  \label{eq:three_wave}
\end{equation}

By tuning the pump, the idler frequency is adjusted:
$\omega_i = 2\omega_p - \omega_s$, enabling continuous frequency
translation of $\psi$ signals.

\subsection{Dual-pump schemes}

Two pumps at $2\omega_1$ and $2\omega_2$ create sum and difference
frequencies:
\begin{equation}
  \omega_{\mathrm{out}} = \omega_s + n(\omega_1 - \omega_2),
  \quad n = \pm 1, \pm 2, \ldots
  \label{eq:dual_pump}
\end{equation}
enabling a comb of $\psi$ channels from a single seed.

\subsection{Application to $\psi$ communications}

Frequency-agile $\psi$ channels enable:
\begin{itemize}[noitemsep]
  \item Frequency-division multiplexing of $\psi$ signals;
  \item Spread-spectrum $\psi$ communication for interference
    rejection;
  \item Frequency hopping for secure $\psi$ links.
\end{itemize}


% ======================================================================
\section{Experimental Roadmap}
\label{sec:roadmap}

\subsection{Phase 1: Proof of concept (Year 1)}

\begin{enumerate}[noitemsep]
  \item Construct LC ladder network with varactor modulation;
  \item Demonstrate parametric EM amplification in the $2\omega$ channel;
  \item Establish noise floor for $\psi$-mode detection;
  \item Target: $|\lambda - 1| \lesssim 10^{-3}$.
\end{enumerate}

\subsection{Phase 2: SRF cavity (Years 2--3)}

\begin{enumerate}[noitemsep]
  \item Commission Nb SRF cavity with TE+TM superposition;
  \item Implement $2\omega$ modulation and PLL control;
  \item Execute full measurement protocol (Sec.~\ref{sec:engineering});
  \item Target: $|\lambda - 1| \lesssim 10^{-8}$.
\end{enumerate}

\subsection{Phase 3: Optimized search (Years 3--5)}

\begin{enumerate}[noitemsep]
  \item Pulsed high-energy operation ($U_0 \to 10^6$~J);
  \item Cavity array with maximized $\Acav$;
  \item Full intentional-search protocol;
  \item Target: $|\lambda - 1| \lesssim 10^{-14}$.
\end{enumerate}

\subsection{Phase 4: Metamaterial arrays (Years 3--5, parallel)}

\begin{enumerate}[noitemsep]
  \item Design and fabricate $8\times 8$ metamaterial panel;
  \item Demonstrate distributed $\psi$ gain and beam-steering;
  \item Scale to $16\times 16$ and beyond;
  \item Integrate with propulsion/energy testbeds.
\end{enumerate}


% ======================================================================
\section{Discussion}
\label{sec:discussion}

\subsection{Why $\lambda \neq 1$ has not been detected}

Three factors explain the null result in existing metrology:
\begin{enumerate}
  \item \textbf{Pure eigenmodes suppress the driven channel.}
    Symmetric cavities in pure modes have $G \approx 0$ by
    Bessel-function orthogonality and equipartition;
  \item \textbf{Parametric pumping needs deliberate $2\omega$.}
    Routine metrology avoids such tones and heavily filters them
    to suppress amplitude-modulation sidebands;
  \item \textbf{$2\omega$ features treated as technical noise.}
    Any residual $2\omega$ response is interpreted as technical AM
    sidebands and actively suppressed, not investigated as a
    potential signal.
\end{enumerate}

The DFD programme does not ask anyone to believe new physics;
it asks them to notice the parametric instability that is not there
and design an experiment specifically to look for it.

\subsection{Implications of detection}

If $|\lambda - 1| \neq 0$ is confirmed:
\begin{enumerate}[noitemsep]
  \item EM fields can generate and amplify $\psi$---the scalar
    refractive field becomes a controllable laboratory resource;
  \item Parametric $\psi$-resonators become practical $\psi$ sources
    for propulsion, energy transfer, and communications;
  \item The DFD framework gains its first direct laboratory evidence
    beyond astronomical observations;
  \item Metamaterial $\psi$-amplifier arrays enable beam-formed,
    frequency-agile $\psi$ technology.
\end{enumerate}

\subsection{Implications of null result}

If $|\lambda - 1| < 10^{-14}$:
\begin{enumerate}[noitemsep]
  \item EM back-reaction on $\psi$ is negligibly small at accessible
    energies---the paddle slides through the water;
  \item Laboratory $\psi$ generation requires alternative coupling
    channels (matter-based, gravitational);
  \item The core DFD framework ($n = e^\psi$, matter coupling,
    field equation) remains valid---$\lambda$ is a distinct
    parameter not constrained by gravitational tests;
  \item Astronomical $\psi$ sources (stellar cores, neutron stars,
    accretion disks) remain operative.
\end{enumerate}


% ======================================================================
\section{Conclusions}
\label{sec:conclusions}

We have developed the complete theoretical framework for parametric
$\psi$-field amplification within Density Field Dynamics, deriving:
\begin{enumerate}[noitemsep]
  \item The parametric gain equation from the DFD action, with
    Mathieu instability threshold and growth rates;
  \item The area-ratio design law connecting cavity geometry to
    the minimum detectable $|\lambda - 1|$;
  \item Small-signal gain curves, bandwidth, stability boundaries,
    and control strategies;
  \item Metamaterial unit-cell designs for distributed $\psi$ gain
    and beam-steering;
  \item Systematic reinterpretation of EMDrive, SRF anomalies,
    JPAs, dynamic Casimir experiments, and time-varying metamaterials
    as potential $\psi$-coupling signatures;
  \item Complete laboratory designs capable of reaching
    $|\lambda - 1| \sim 10^{-14}$.
\end{enumerate}

The central message is operational: existing high-$Q$ cavities
\emph{accidentally} constrain $|\lambda - 1| \lesssim 3\times 10^{-5}$
through the absence of parametric instability.  A deliberate search,
using the same hardware with TE+TM superposition and preserved $2\omega$
response, pushes this bound ten orders of magnitude deeper.  A single
afternoon's measurement could either discover $\lambda \neq 1$ or
constrain it below $10^{-14}$.

% ======================================================================
\begin{acknowledgments}
The author thanks the accelerator physics and superconducting RF
communities for decades of high-$Q$ cavity data that inadvertently
constrain the EM--$\psi$ back-reaction parameter.
\end{acknowledgments}

\begin{thebibliography}{99}

\bibitem{Alcock2025DFD}
G.~T.~Alcock,
``Density Field Dynamics: A Complete Unified Theory,''
v3.2 (March 2026).

\bibitem{Alcock2025Patent2}
G.~T.~Alcock,
``Parametric $\psi$-Resonators and Metamaterial Amplifiers,''
U.S. Provisional Patent Application (October 11, 2025).

\bibitem{Alcock2025EMCoupling}
G.~T.~Alcock,
``Systems and Methods for Electromagnetic Coupling, Actuation, and
Application of Scalar Refractive Fields ($\psi$),''
U.S. Provisional App.\ No.\ 63/893,633 (October 4, 2025).

\bibitem{Shawyer2006}
R.~Shawyer,
``A Theory of Microwave Propulsion for Spacecraft,''
New Scientist \textbf{2568}, 30 (2006).

\bibitem{White2017}
H.~White, P.~March, J.~Lawrence, J.~Vera, A.~Sylvester, D.~Brady,
and P.~Bailey,
``Measurement of Impulsive Thrust from a Closed Radio-Frequency Cavity
in Vacuum,''
J.\ Propuls.\ Power \textbf{33}, 830 (2017).

\bibitem{Tajmar2021}
M.~Tajmar, M.~Kossling, M.~Weikert, and M.~Monette,
``The SpaceDrive Project---Thrust Balance Development and New
Measurements of the Mach-Effect and EMDrive Thrusters,''
Acta Astronaut.\ \textbf{187}, 224 (2021).

\bibitem{Wilson2011}
C.~M.~Wilson, G.~Johansson, A.~Pourkabirian, M.~Simoen,
J.~R.~Johansson, T.~Duty, F.~Nori, and P.~Delsing,
``Observation of the dynamical Casimir effect in a superconducting
circuit,''
Nature \textbf{479}, 376 (2011).

\bibitem{Lahteenmaki2013}
P.~L\"ahteenm\"aki, G.~S.~Paraoanu, J.~Hassel, and P.~J.~Hakonen,
``Dynamical Casimir effect in a Josephson metamaterial,''
Proc.\ Natl.\ Acad.\ Sci.\ \textbf{110}, 4234 (2013).

\bibitem{Galiffi2022}
E.~Galiffi, R.~Tirole, S.~Yin, H.~Li, S.~Vezzoli, P.~A.~Huidobro,
M.~G.~Silveirinha, R.~Sapienza, A.~Al\`u, and J.~B.~Pendry,
``Photonics of time-varying media,''
Adv.\ Photon.\ \textbf{4}, 014002 (2022).

\bibitem{Castellanos2008}
A.~Castellanos-Beltran, K.~Irwin, G.~Hilton, L.~Vale,
and K.~Lehnert,
``Amplification and squeezing of quantum noise with a tunable
Josephson metamaterial,''
Nat.\ Phys.\ \textbf{4}, 929 (2008).

\bibitem{Padilla2022}
W.~J.~Padilla and R.~D.~Averitt,
``Properties of dynamical electromagnetic metamaterials,''
J.\ Opt.\ Soc.\ Am.\ B \textbf{39}, 1514 (2022).

\bibitem{Liu2000}
Z.~Liu, X.~Zhang, Y.~Mao, Y.~Y.~Zhu, Z.~Yang, C.~T.~Chan,
and P.~Sheng,
``Locally Resonant Sonic Materials,''
Science \textbf{289}, 1734 (2000).

\bibitem{Casimir1948}
H.~B.~G.~Casimir,
``On the attraction between two perfectly conducting plates,''
Proc.\ K.\ Ned.\ Akad.\ Wet.\ \textbf{51}, 793 (1948).

\bibitem{Padilla2015}
W.~J.~Padilla, A.~J.~Taylor, C.~Highstrete, M.~Lee,
and R.~D.~Averitt,
``Dynamical Electric and Magnetic Metamaterial Response at Terahertz
Frequencies,''
Phys.\ Rev.\ Lett.\ \textbf{96}, 107401 (2006).

\end{thebibliography}

\end{document}
